%% The full diagram gives a CCC equivalence with explicit comparison maps.

\section{The internal language of a small cartesian closed category}

\begin{annexitem}{annex:p10.internal.1}{thm:p10-internal-ccc-comparisons}{Type comparison and context packing}
\begin{proof}
\paragraph{Equality in the diagram theory.}
The typing and substitution inductions of
Theorem~\ref{thm:p10-stlc-equational-substitution} extend to closed
constants: a constant remains unchanged by substitution and retains
its type under weakening and renaming.
In a context, its denotation is its closed denotation composed with
the unique arrow to \(1\).  Thus the constant case of the semantic
substitution law is the identity \(!\circ a=!\).
The other induction cases are those of
Theorem~\ref{thm:p10-stlc-ccc-completeness}.

Every generator of \(=_{\mathcal C}\) is sound in the given
interpretation.  This is the definition for diagram axioms and is
the categorical \(\beta\eta\) law for those equations.
Equivalence and congruence preserve semantic equality; abstraction
uses transposition, and substitution uses composition.
Consequently
\[
 \Theta_\Gamma\vdash M=_{\mathcal C}N:T
 \quad\Longleftrightarrow\quad
 \llbracket M\rrbracket_\Gamma=\llbracket N\rrbracket_\Gamma.
\]
For the reverse implication the equation is itself a diagram axiom.
This equivalence is a property of the explicitly defined full diagram.

\paragraph{Comparing a type with its named object.}
For \(T=\underline A\), let both \(\kappa_T\) and \(\iota_T\)
be the identity abstraction.
Suppose comparisons have been defined for \(S,T\), and put
\(A=A_S\), \(B=A_T\).  Define the closed terms
\[
\begin{aligned}
 \kappa_{S\to_{\rm ty}T}
  &=\lambda h{:}(S\to_{\rm ty}T).
      l_{A,B}\bigl(\lambda a{:}\underline A.
                        \kappa_T(h(\iota_Sa))\bigr),\\
 \iota_{S\to_{\rm ty}T}
  &=\lambda z{:}\underline{B^A}.\lambda s{:}S.
                   \iota_T\bigl((e_{A,B}z)(\kappa_Ss)\bigr).
\end{aligned}
\]
All the displayed applications have the required types.
By induction the inner comparisons act as identities on \(A,B\).
More precisely, for any parameter object \(D\) and arrow
\(r:D\to B^A\) interpreting the variable \(h\), the first inner
body denotes \(\operatorname{ev}_{A,B}\circ(r\times1_A)\).
Its abstraction denotes
\[
 \Lambda\bigl(\operatorname{ev}_{A,B}\circ(r\times1_A)\bigr)=r
\]
by the exponential \(\eta\) identity.
The same calculation applies to the second abstraction, with the
projection interpreting \(z\) in place of \(r\).
Since \(e_{A,B}\) and \(l_{A,B}\) have identity unary interpretations,
both new comparisons act as the identity on \(B^A\).
Equivalently, applying either comparison to a variable denotes the
corresponding variable projection in its singleton context.
Their composites therefore have the identity denotations, so their
two inverse equations belong to the diagram theory.

\paragraph{Packing and unpacking the finite context.}
For a context with chosen variable names, define a packing term
\(p_\Gamma:\underline{P_\Gamma}\) recursively by
\[
 p_{()}=u,\qquad
 p_{\Gamma\cdot(x{:}T)}
   =\mathsf{pair}_{P_\Gamma,A_T}\,p_\Gamma\,(\kappa_Tx),
\]
weakening \(p_\Gamma\) into the extended context.
The empty term denotes \(1_1\).
At the successor step the weakened prefix denotes \(\pi_1\), the
last comparison denotes \(\pi_2\), and pairing gives
\(\langle\pi_1,\pi_2\rangle=1_{P_\Gamma\times A_T}\).
Thus \(\llbracket p_\Gamma\rrbracket=1_{P_\Gamma}\) for every context.

Let \(\operatorname{pack}_\Gamma\) be the one-component tuple
\((p_\Gamma)\).
Its target context is \((\underline{P_\Gamma})\), whose interpretation
is \(1\times P_\Gamma\).  Its context-arrow interpretation is
therefore \(s_{P_\Gamma}=\langle!,1_{P_\Gamma}\rangle\).
For \(\Gamma=(T_1,\ldots,T_m)\), let
\(\rho_i:P_\Gamma\to A_{T_i}\) be its variable projections and define
\(\operatorname{unpack}_\Gamma\) in context
\(z:\underline{P_\Gamma}\) by the tuple
\[
 \bigl(\iota_{T_i}(c_{\rho_i}z)\bigr)_{i=1}^m.
\]
Its components denote \(\rho_i\circ\pi_2\).
The finite-product universal property, including its unique terminal
component, shows that the tuple denotes
\(\varepsilon_{P_\Gamma}=\pi_2:1\times P_\Gamma\to P_\Gamma\).

Here for every object \(A\) put
\(s_A=\langle!_A,1_A\rangle\) and \(\varepsilon_A=\pi_2\).
Terminal uniqueness and the product laws give
\[
 \varepsilon_A\circ s_A=1_A,\qquad
 s_A\circ\varepsilon_A=1_{1\times A}.
\]
The composite packing/unpacking tuples consequently have identity
denotations, and their component equations are in the diagram.
If \(\Gamma=()\), unpacking is the empty tuple and its denotation
is the unique map \(1\times1\to1\).
The other composite is the equation \(z=u\) in type \(\underline1\),
which is sound by terminal uniqueness.  This proves both inverse
equations in every case.
\end{proof}
\end{annexitem}

\begin{annexitem}{annex:p10.internal.2}{thm:p10-internal-ccc-equivalence}{Full faithfulness and an explicit quasi-inverse}
\begin{proof}
\paragraph{Evaluation is a functor.}
The category laws for typed tuples follow from substitution identities,
as in Definition~\ref{def:p10-stlc-context-category}.
The added constants and diagram equations preserve the required
substitution and congruence properties by the preceding proof.
For a tuple \(M:\Gamma\to\Delta\), define \(E(M):P_\Gamma\to P_\Delta\)
to be the unique arrow with its denoted term components.
The leading terminal component is the unique arrow to \(1\).
Semantic soundness makes this independent of the representative tuple.
The substitution lemma proves composition preservation, and variable
projections prove preservation of identities.

\paragraph{Fullness.}
Given \(f:P_\Gamma\to P_\Delta\), use the context arrow
\[
 M_f=\operatorname{unpack}_\Delta\circ(c_fx)
                         \circ\operatorname{pack}_\Gamma.
\]
The middle arrow has source \((\underline{P_\Gamma})\) and target
\((\underline{P_\Delta})\).  Its interpretation is
\(1_1\times f:1\times P_\Gamma\to1\times P_\Delta\).
Consequently
\[
 E(M_f)=\varepsilon_{P_\Delta}\circ(1_1\times f)\circ s_{P_\Gamma}=f.
\]
The equality follows by the product projections.
For an empty source, the packing term is \(u\); for an empty target,
the unpacking tuple is empty and \(f\) is the unique map to \(1\).
Thus these cases also establish fullness.

\paragraph{Faithfulness and essential surjectivity.}
If \(E(M)=E(N)\), compose with each variable projection of the target
context.  Corresponding component terms then have equal denotations,
so their equations are diagram axioms.  Hence the tuples give the
same arrow in \(\mathcal S_{\mathcal C}\).
For an empty target there is already only one tuple.
This proves faithfulness.
For each \(A\in\mathcal C\), the context \((\underline A)\) evaluates
to \(1\times A\), isomorphic to \(A\) through \(\varepsilon_A\).
Thus \(E\) is essentially surjective with an explicit representative.

\paragraph{Quasi-inverse, naturality, and triangles.}
Define \(J(A)=(\underline A)\) and \(J(f)=(c_fx)\).
The identity and composition equations in the diagram make \(J\) a
functor.  Its composite \(EJ\) sends \(A\) to \(1\times A\) and
\(f\) to \(1_1\times f\).
The projection identity
\[
 \varepsilon_B\circ(1_1\times f)=f\circ\varepsilon_A
\]
proves that \(\varepsilon:EJ\Rightarrow1_{\mathcal C}\) is a natural
isomorphism, with components inverse to \(s_A\).

Put \(\eta_\Gamma=\operatorname{pack}_\Gamma:\Gamma\to JE(\Gamma)\).
For \(M:\Gamma\to\Delta\), the arrows
\(JE(M)\circ\eta_\Gamma\) and \(\eta_\Delta\circ M\) both evaluate to
\(s_{P_\Delta}\circ E(M)\), because
\((1_1\times E(M))\circ s_{P_\Gamma}=s_{P_\Delta}\circ E(M)\).
Faithfulness proves naturality of \(\eta\); its inverse components
are the unpacking arrows already constructed.
The first triangle is
\[
 \varepsilon_{E\Gamma}\circ E(\eta_\Gamma)
     =\varepsilon_{P_\Gamma}\circ s_{P_\Gamma}=1_{E\Gamma}.
\]
For the second, in context \(x:\underline A\), the composite
\(J(\varepsilon_A)\circ\eta_{J(A)}\) has the term
\[
 c_{\pi_2}\bigl(\mathsf{pair}_{1,A}\,u\,x\bigr)=x.
\]
Here the base-type comparison is the identity term, and the displayed
equation is the product projection equation in the diagram.
Thus the second triangle is \(1_{J(A)}\).
These explicit functors and natural isomorphisms prove the equivalence
without invoking a choice of representatives from arbitrary isomorphism classes.
\end{proof}
\end{annexitem}

\begin{annexitem}{annex:p10.internal.3}{thm:p10-internal-ccc-structure}{The cartesian closed comparison isomorphisms}
\begin{proof}
\paragraph{Cartesian closure of the diagram category.}
The terminal and product proofs of
Theorem~\ref{thm:p10-stlc-context-ccc} use only empty tuples,
concatenation, and substitution, and remain valid here.
For \(\Gamma=(S_1,\ldots,S_m)\), \(\Delta=(T_1,\ldots,T_n)\),
take \(\Omega=(c_\Gamma(T_1),\ldots,c_\Gamma(T_n))\).
Application and abstraction over the \(m\) variables give evaluation
and transposition.  Congruence makes these operations well-defined
modulo the diagram equations.  Their two inverse laws follow from
successive \(\beta\) and \(\eta\) equations, and substitution gives
naturality in the parameter context, exactly as in that theorem.
Thus the larger signature and equation set retain cartesian closure.

\paragraph{Terminal and product comparisons.}
The empty context has interpretation \(1\).
The arrow \(j_{\Gamma,\Delta}\) in the statement has as components
the interpreted block projections.
Its inverse is the unique arrow into \(P_{\Gamma\cdot\Delta}\)
whose variable components first project from \(P_\Gamma\), then from
\(P_\Delta\).  Both composites are identities because they have the
same variable projections as the identities; the terminal components
are unique.  This argument includes empty blocks.
The same projection calculation gives naturality:
\[
 j_{\Gamma',\Delta'}\circ E(M\times N)
  =(E(M)\times E(N))\circ j_{\Gamma,\Delta}
\]
for \(M:\Gamma\to\Gamma'\), \(N:\Delta\to\Delta'\).
All unit and associativity comparisons obtained this way commute,
since their parallel composites have identical product projections.

\paragraph{The interpreted exponential has its universal property.}
Put \(A=P_\Gamma\), \(B=P_\Delta\), and
\[
 v=E(\operatorname{ev}_{\rm syn})\circ j_{\Omega,\Gamma}^{-1}
                      :P_\Omega\times A\to B.
\]
Let \(D\) be any object and \(f:D\times A\to B\).
Full faithfulness gives a unique syntactic arrow
\(M:J(D)\cdot\Gamma\to\Delta\) whose interpretation is
\[
 E(M)=f\circ(\varepsilon_D\times1_A)\circ j_{J(D),\Gamma}.
\]
Transpose \(M\) syntactically to \(\widetilde M:J(D)\to\Omega\),
and put
\[
 q=E(\widetilde M)\circ s_D:D\to P_\Omega.
\]
Naturality of \(j\), followed by the syntactic evaluation law, gives
\[
 v\circ(E(\widetilde M)\times1_A)
       =E(M)\circ j_{J(D),\Gamma}^{-1}
       =f\circ(\varepsilon_D\times1_A).
\]
Composing with \(s_D\times1_A\) yields \(v\circ(q\times1_A)=f\).

For uniqueness, suppose \(q':D\to P_\Omega\) has that property.
Full faithfulness gives a unique \(Q:J(D)\to\Omega\) with
\(E(Q)=q'\circ\varepsilon_D\).
The same product calculation shows that the syntactic uncurrying
of \(Q\) has interpretation \(E(M)\).  Faithfulness identifies it
with \(M\), and syntactic transpose uniqueness gives
\(Q=\widetilde M\).  Composing their interpretations with \(s_D\)
gives \(q'=q\).
Thus \((P_\Omega,v)\) is an exponential of \(B\) by \(A\).

The arrow \(\chi=\Lambda(v):P_\Omega\to B^A\) is its comparison
with the given exponential.  The just-proved universal property
gives a unique \(r:B^A\to P_\Omega\) with
\(v\circ(r\times1_A)=\operatorname{ev}_{A,B}\).
The two exponential uniqueness clauses show
\(\chi r=1_{B^A}\) and \(r\chi=1_{P_\Omega}\).
Hence \(\chi\) is an isomorphism and is precisely the map in the statement.

If \(\Gamma=()\), then \(\Omega=\Delta\) and this compares
\(P_\Delta\) with \(P_\Delta^1\).
If \(\Delta=()\), then \(\Omega=()\) and it compares \(1\) with
\(1^{P_\Gamma}\); the same uniqueness proofs apply.
For the quasi-inverse, the constants \(u\), \(\mathsf{pair}\), and
\(e,l\) give respectively the terminal, product, and exponential
comparisons for named objects.  The structural diagram equations
ensure agreement with projections and evaluation.
This proves preservation of cartesian closed structure up to the
specified canonical comparison isomorphisms.
\end{proof}
\end{annexitem}
